% # Generate noise images 
% from PIL import Image
% import torch
% n = 3
% size = 16
% levels = [0.25, 1.0]
% for i in range(n):
%     for level in levels:
%         x = torch.randn(size, size, 3).div(1.5).clip(-1,1).add(1).div(2)
%         m = torch.rand(size, size, 1) > level
%         m = torch.cat([m] * 3, dim=-1)
%         x[m] = 1
%         Image.fromarray(x.mul(255).byte().numpy()).save(f"noise_{i}_level_{level}.png")

\begin{SCfigure}[][h]
%\begin{wrapfigure}{R}{0.5\linewidth}
% \begin{figure}[t!]

\vspace{-5mm}
\resizebox{0.5\linewidth}{!}{%

\begin{tikzpicture}[
    image/.style={black,fill=white,draw,inner sep=0,minimum width=1.1cm,minimum height=1.1cm},
    maps/.style={->,shorten <=1pt,shorten >=1pt,}
  ]

\def\horigap{0.75cm}
\def\vertgap{0.75cm}
\def\smallvertgap{-0.15cm}

%%%%%%%%%%%%%%%%%%%%%%%%%%%%%%%%%%%%%%%%%%%%
% t=0

\path
node[image] (cond n=0 t=0) at (0, 0) {}
node[image] (cond n=1 t=0) at ($(\tikzlastnode)+(3pt,-2pt)$) {}
node[image] (cond n=2 t=0) at ($(\tikzlastnode)+(3pt,-2pt)$) {}
node[image] (cond n=3 t=0) at ($(\tikzlastnode)+(3pt,-2pt)$) {}
node[image] (cond n=4 t=0) at ($(\tikzlastnode)+(3pt,-2pt)$) {$\x^0_{t-1}$}
node[image,above=\vertgap of \tikzlastnode,xshift=-6pt] (X k=0 t=0) {\includegraphics[width=1.1cm]{images/noise/noise_0_level_1.0.png}}
node[minimum height=1cm,above=\vertgap of \tikzlastnode] (etc 1 t=0) {$\vdots$}
node[image,above=\smallvertgap of \tikzlastnode] (X k=1 t=0) {\includegraphics[width=1.1cm]{images/noise/noise_0_level_0.5.png}}
node[minimum height=1cm,above=\vertgap of \tikzlastnode] (etc 2 t=0) {$\vdots$}
node[image,above=\smallvertgap of \tikzlastnode] (X k=K t=0) {}
node[right=\horigap of \tikzlastnode,inner sep=2pt,circle,minimum width=1.1cm,draw] (A t=0) {$a_t$}
;

\node[fill=white,inner sep=2pt] at (X k=0 t=0) {$\x^\Tau_t$};
\node[draw=none,fill=white,inner sep=2pt] at (X k=1 t=0) {$\x^\tau_t$};
\node[draw=none,fill=white,inner sep=2pt] at (X k=K t=0) {$\x^0_t$};

%%%%%%%%%%%%%%%%%%%%%%%%%%%%%%%%%%%%%%%%%%%%
% t=1

\path
node[image,right=3cm of {cond n=0 t=0}] (cond n=0 t=1) at (0, 0) {}
node[image] (cond n=1 t=1) at ($(\tikzlastnode)+(3pt,-2pt)$) {}
node[image] (cond n=2 t=1) at ($(\tikzlastnode)+(3pt,-2pt)$) {}
node[image] (cond n=3 t=1) at ($(\tikzlastnode)+(3pt,-2pt)$) {}
node[image] (cond n=4 t=1) at ($(\tikzlastnode)+(3pt,-2pt)$) {$\x^0_t$}
node[image,above=\vertgap of \tikzlastnode,xshift=-6pt] (X k=0 t=1) {\includegraphics[width=1.1cm]{images/noise/noise_1_level_1.0.png}}
node[minimum height=1cm,above=\vertgap of \tikzlastnode] (etc 1 t=1) {$\vdots$}
node[image,above=\smallvertgap of \tikzlastnode] (X k=1 t=1) {\includegraphics[width=1.1cm]{images/noise/noise_1_level_0.5.png}}
node[minimum height=1cm,above=\vertgap of \tikzlastnode] (etc 2 t=1) {$\vdots$}
node[image,above=\smallvertgap of \tikzlastnode] (X k=K t=1) {}
node[right=\horigap of \tikzlastnode,inner sep=2pt,circle,minimum width=1.1cm,draw] (A t=1) {$a_{t+1}$}
;

\node[fill=white,inner sep=2pt] at (X k=0 t=1) {$\x^\Tau_{t+1}$};
\node[draw=none,fill=white,inner sep=2pt] at (X k=1 t=1) {$\x^\tau_{t+1}$};
\node[draw=none,fill=white,inner sep=2pt] at (X k=K t=1) {$\x^0_{t+1}$};

%%%%%%%%%%%%%%%%%%%%%%%%%%%%%%%%%%%%%%%%%%%%
% t=T

\path
node[image,right=7cm of {cond n=0 t=1}] (cond n=0 t=T) at (0, 0) {}
node[image] (cond n=1 t=T) at ($(\tikzlastnode)+(3pt,-2pt)$) {}
node[image] (cond n=2 t=T) at ($(\tikzlastnode)+(3pt,-2pt)$) {}
node[image] (cond n=3 t=T) at ($(\tikzlastnode)+(3pt,-2pt)$) {}
node[image] (cond n=4 t=T) at ($(\tikzlastnode)+(3pt,-2pt)$) {$\x^0_{T-1}$}
node[image,above=\vertgap of \tikzlastnode, xshift=-6pt] (X k=0 t=T) {\includegraphics[width=1.1cm]{images/noise/noise_2_level_1.0.png}}
node[minimum height=1cm,above=\vertgap of \tikzlastnode] (etc 1 t=T) {$\vdots$}
node[image,above=\smallvertgap of \tikzlastnode] (X k=1 t=T) {\includegraphics[width=1.1cm]{images/noise/noise_2_level_0.5.png}}
node[minimum height=1cm,above=\vertgap of \tikzlastnode] (etc 2 t=T) {$\vdots$}
node[image,above=\smallvertgap of \tikzlastnode] (X k=K t=T) {}
%% node[right=\horigap of \tikzlastnode,inner sep=2pt,circle,minimum width=1.1cm,draw] (A t=T) {$$}
;

\node[fill=white,inner sep=2pt] at (X k=0 t=T) {$\x^\Tau_{T}$};
\node[draw=none,fill=white,inner sep=2pt] at (X k=1 t=T) {$\x^\tau_{T}$};
\node[draw=none,fill=white,inner sep=2pt] at (X k=K t=T) {$\x^0_{T}$};

%%%%%%%%%%%%%%%%%%%%%%%%%%%%%%%%%%%%%%%%%%%%%%%%%%%%%%%%%%%%%%%%%%%%%%

\foreach \t in {0,1}{
  \draw[maps] (X k=K t=\t)--(A t=\t) node[midway,above] {$\pi_\phi$};
}

\foreach \t in {0,1,T}{
%%   \foreach \a/\b in {X k=0 t=\t/etc 1 t=\t, etc 1 t=\t/X k=1 t=\t,X k=1 t=\t/etc 2 t=\t,etc 2 t=\t/X k=K t=\t}{
  \foreach \a/\b in {X k=0 t=\t/etc 1 t=\t,X k=1 t=\t/etc 2 t=\t}{
    \draw[maps] (\a) -- (\b) node[midway,right] {$\mathbf{D}_\theta$};
  }
}

\begin{scope}[on background layer]
  \foreach \n/\a in {n=0 t=0/A,n=1 t=0/A,n=2 t=0/A,n=3 t=0/A,n=4 t=0/$a_{t-1}$,n=0 t=1/A,n=1 t=1/A,n=2 t=1/A,n=3 t=1/A,n=4 t=1/$a_t$,n=0 t=T/A,n=1 t=T/A,n=2 t=T/A,n=3 t=T/A,n=4 t=T/$a_{T-1}$}{
    \node[below=4pt of {cond \n},fill=white,inner sep=2pt,circle,minimum width=1.1cm,draw] (A cond \n) {\a};
  };
\end{scope}

% \node (temp etc) at ($(X k=0 t=1)!0.5!(X k=0 t=T)$) {$\dots$};

\node (temp etc) at ($(X k=K t=1)!0.73!(X k=K t=T)$) {$\dots$};

\draw[->,very thick] ($(X k=K t=0.north west)+(0,5mm)$)--($(X k=K t=T.north east)+(0,5mm)$) node[midway,above,inner sep=5pt] {Environment time ($t$)};

\draw[->,very thick] ($(X k=0 t=0.south west)+(-5mm,0)$)--($(X k=K t=0.north west)+(-5mm,0)$) node[midway,above,inner sep=5pt,rotate=90] (denoising) {Denoising time ($\tau$)};

\draw[draw=none,very thick] ($(X k=0 t=0.south west)+(-5mm,0)$)--++(0,-3) node[midway,above,inner sep=5pt,rotate=90] (conditioning) {Conditioning};

%\node[left=20pt of {cond n=0 t=0},rotate=90,anchor=south] (label) {Conditioning};

%% \node[below=5cm of denoising,rotate=90,yshift=5.5mm] (label) {Conditioning};

%% \draw[->,blue!50,very thick,shorten <=1pt,shorten >=1pt,preaction={draw=white,line width=2pt,-}] (X k=K t=0.south east) to[out=335,in=150] (cond n=4 t=1);

%%%%%%%%%%%%%%%%%%%%%%%%%%%%%%%%%%%%%%%%%%%%

\end{tikzpicture}
}

\caption{Unrolling imagination of \textsc{diamond} over time. The top row depicts a policy $\pi_\phi$ taking a sequence of actions in the imagination of our learned diffusion world model $\mathbf{D}_\theta$. The environment time $t$ flows along the horizontal axis, while the vertical axis represents the denoising time $\tau$ flowing backward from $\Tau$ to $0$. Concretely, given (clean) past observations $\x^0_{<t}$, actions $a_{<t}$, and starting from an initial noisy sample $\x_t^\Tau$, we simulate a reverse noising process $\{\x_t^\tau\}_{\tau =\Tau}^{\tau=0}$ by repeatedly calling $\mathbf{D}_\theta$, and obtain the (clean) next observation $\x_t^0$. The imagination procedure is autoregressive in that the predicted observation $\x_t^0$ and the action $a_t$ taken by the policy become part of the conditioning for the next time step. Animated visualizations of this procedure can be found at \url{https://diamond-wm.github.io}.}

\end{SCfigure}
\label{fig:architecture}